\documentclass{article}
\usepackage{geometry}
\geometry{a4paper, margin=1in}
\title{ML Research Proposal: Causal Inference for Deciphering Morphogenetic Network Topology}
\author{Candidate Name}
\date{Suggested Supervisors: Dr. Carlo Vittorio Cannistraci, Dr. Elias Barriga}

\begin{document}
\maketitle

\section*{Abstract}
Morphogenesis is controlled by a complex regulatory network where bioelectric signals, mechanical forces, and gene expression are tightly coupled, but the direction of causality is often unknown (e.g., does voltage change cause force, or vice-versa?). Deciphering this network topology is crucial for therapeutic reprogramming, aligning with the "compiler" approach. This project uses \textbf{Machine Learning for Causal Inference} to establish the information flow within this biological network. We will analyze multi-modal, time-series data from bioelectric/mechanical assays (e.g., $\text{V}_{\text{mem}}$ reporters, FRET-based tension sensors, and gene expression reporters). Techniques like \textbf{Dynamic Causal Modeling (DCM)} or \textbf{Granger Causality} will be employed to determine statistical dependencies and time-lagged influence between these signals. The main goal is to generate a probabilistic, weighted, and directed graph representing the regulatory logic of the system, ultimately identifying the true upstream control nodes (the "master switches") in the bio-mechanical network. This provides a data-driven complement to the mechanistic simulations of TopoSPAM \cite{1}.

\section*{Related Work}
\begin{itemize}
    \item \textbf{Network Science/Information Theory in Biomedicine:} Application of network-based ML methods to infer regulatory logic and connectivity in complex biological systems (Cannistraci Group expertise).
    \item \textbf{Bioelectric Signalling as a Developmental Cue:} The role of membrane voltage as a high-level signaling component that coordinates collective cell behavior (Levin, 2012; Barriga Lab data).
    \item \textbf{Multi-Omics Time Series Analysis:} Use of temporal data to build dynamic models of biological networks.
\end{itemize}

\end{document}
